\section{Households, by position in the wage distribution}
\label{sec:households}

The accounts carry one more cut. Because the household part of the wage-backed stock can be split
by wage quintile, and because the same survey gives the liquid assets of the same households, the
two can be set side by side. They answer different questions and they disagree, which is the
result of this section.

\subsection{Claims are spread more evenly than the income that services them}

The gradient is the finding: wage-backed claims per dollar of wages fall as pay rises, from
\result{QTwoPerDollar} in the second quintile to \result{QTopPerDollar} at the top, while the top
quintile earns \result{QTopWageShare} percent of the wage bill and carries
\result{QTopClaimShare} percent of the household claims. Claims are spread far more evenly than
the income that services them. Table~\ref{tab:quintile} gives the detail, and
Appendix~\ref{app:detail} the construction, including the one magnitude the accounts can produce
but that we do not report.

That is the sense in which a proportional fall in pay is regressive before any policy response is
chosen. It is an arithmetic property of the claim structure, not a claim about how anyone
behaves. A household in the second quintile supplies about nine percent of the wage bill and
carries about thirteen percent of the household claims against it; a household in the top
quintile supplies half the wage bill and carries a third of the claims. Nothing has to happen for
that asymmetry to exist. It is there now.

\subsection{Where the buffers are thinnest, and it is not the bottom}

Liquid runway by wage quintile gives the second household result, and buffers are not monotonic
in pay: the thinnest are in the \emph{second} quintile, a median of \result{QTwoRunway} months
against \result{QOneRunway} in the bottom quintile, with \result{QTwoUnderOne} percent under one
month against \result{QOneUnderOne} percent. From there upward the gradient is monotonic,
reaching \result{QFiveRunway} months at the top. Bottom-quintile households hold little debt and
few of the commitments a runway is measured against, while second-quintile households carry the
commitments of those above them without the income to service them, so relief targeted at the
bottom of the distribution misses the households with the least room.

\begin{table}[htbp]
\centering
\caption{Liquid runway of working households by wage quintile: months of income covered by liquid
assets, and the share with less than one month. Measured.}
\label{tab:buffers}
\input{tables/tab_buffers.tex}
\end{table}

\subsection{What the two cuts say together}

The claims gradient and the buffer gradient are measured on the same households from the same
survey, and they do not line up. Claims per wage dollar peak one quintile above the bottom.
Liquid runway bottoms out in the same place. The households carrying the most debt per dollar of
income earned are also the households with the least room to absorb a gap in that income, and
neither fact is visible in a statistic that ranks households by income alone or by wealth alone.

We state this as a measured coincidence rather than as a mechanism. The accounts cannot say why
the second quintile occupies that position, and a plausible reading, that these are households
with enough income to qualify for credit and not enough to build a buffer against it, is a
hypothesis this paper does not test. What the accounts establish is the location. Anything
targeted at the bottom of the wage distribution, whether credit policy, relief or disclosure,
is aimed one quintile away from the thinnest margin in the household sector.
